\chapter{Motivación}


\section{El desafío de prevenir la ceguera irreversible en personas con diabetes}
La diabetes es una enfermedad metabólica crónica caracterizada por altos niveles de azúcar en sangre, que con el tiempo puede causar daños graves en diversos órganos, como el corazón, los riñones, los ojos, así como en los vasos sanguíneos y los nervios \parencites{ops_diabetes}. En el mundo, más de 537 millones de adultos, entre 20 y 79 años, viven actualmente con diabetes. Se estima que esta cifra aumentará a 783 millones para 2045. Además, los gastos asociados con la enfermedad se han incrementado en un 316\% en los últimos 15 años, alcanzando al menos 966 mil millones de dólares \parencites{idf_2021}.


En Argentina, se registró un aumento significativo de la prevalencia de glucemia elevada/diabetes en la población total (18 años y más), pasando del 9,8\% en 2013 al 12,7\% en 2018. Además, el 19,8\% de la población presenta riesgo alto o muy alto de desarrollar diabetes mellitus en los próximos diez años  \parencites{indec2019}. 

Existen ciertas complicaciones oculares relacionadas con la diabetes, como la retinopatía diabética (RD), que se produce cuando los elevados niveles de glucosa dañan los vasos sanguíneos de la retina. Con el tiempo, los vasos sanguíneos dañados pueden perder líquido o sangrar provocando alteraciones en la vista. A nivel global, se estima que aproximadamente un tercio de las personas con diabetes desarrolla algún grado de RD, siendo esta la principal causa de ceguera en adultos en edad laboral \parencites{idf_2021} y la segunda causa de ceguera en Argentina (16\%) \parencites{silva2015}.

Al ser una enfermedad crónica, es fundamental instruir a la población sobre el autocuidado adecuado y el seguimiento de la enfermedad y sus complicaciones. Por ello, las guías recomiendan realizar anualmente una imagen de fondo de ojo \parencites{ops_diabetes, ops_clinica_diabetes}. La detección temprana, junto con el tratamiento oportuno y un seguimiento adecuado, puede reducir hasta en un 95\% el riesgo de pérdida de visión asociada con la RD \parencites{nei2023}.

En una zona rural de La Pampa, un estudio reveló que más del 30\% de las personas evaluadas nunca se había realizado un fondo de ojos. Además, el riesgo de desarrollar RD era significativamente mayor en quienes no se sometían a este examen desde hace más de cinco años, destacando la falta de adherencia a esta práctica esencial para la detección temprana \parencites{ortizbasso2022}.

En relación a la prevalencia de la RD en Argentina, la Campaña de Prevención de Ceguera por Diabetes, en el año 2023, identificó un 29.4\% de RD \parencites{oftalmologos2023}.


\section{Barreras en el diagnóstico temprano de la retinopatía diabética y sus consecuencias}
En Argentina, las desigualdades en el acceso al diagnóstico y tratamiento de la diabetes y sus complicaciones, como la retinopatía diabética, están marcadas por diferencias socioeconómicas y geográficas. Según la Encuesta Nacional de Factores de Riesgo, la prevalencia de diabetes por autorreporte varía considerablemente entre regiones: en la Ciudad de Buenos Aires es del 8,8\%, mientras que en la provincia de San Luis alcanza el 17,3\%, siendo esta la más alta del país.

El nivel educativo también juega un rol importante, ya que las personas con menor nivel educativo (hasta primario incompleto) tienen una prevalencia de diabetes de 22,9\%, en contraposición al 10,3\% de quienes alcanzaron el nivel secundario completo o superior. Estas diferencias reflejan la influencia del acceso a la información y los recursos preventivos.

En relación a los estudios para la detección temprana de complicaciones, las desigualdades también son evidentes. Las personas de sexo femenino reportan menos exámenes de fondo de ojo que el sexo masculino, y las personas que dependen exclusivamente del sistema público de salud realizan menos estudios de la vista en comparación con aquellas que cuentan con algún tipo de cobertura médica privada. Asimismo, el ingreso económico es un factor relevante: mientras que solo el 30,1\% de las personas en el quintil de ingresos más bajos realizó estudios de la vista, este porcentaje asciende al 47,8\% en el quintil más alto \parencites{indec2019}.

Por otro lado, en las zonas rurales de Argentina, las barreras para acceder al diagnóstico temprano se ven exacerbadas por la centralización del sistema de salud, que obliga a las personas a recorrer largas distancias a los centros de salud ya sea para un exámen o recibir el tratamiento correspondiente \parencites{ortizbasso2020}.

Cuando la RD se detecta en etapas tempranas, seguir una dieta saludable, controlar los niveles de azúcar en sangre y la presión arterial puede prevenir su progresión a formas más graves e incluso mejorar parcialmente la visión. Sin embargo, si la RD no se diagnostica hasta etapas más avanzadas, el tratamiento necesario es más complejo y costoso, e incluso puede requerir cirugía para reparar el daño en la retina. El diagnóstico temprano no solo mejora los resultados de salud, sino que también tiene un impacto económico positivo tanto para los pacientes como para el sistema de salud \parencites{aao2023}.


\section{Potencial de las soluciones basadas en inteligencia artificial}
En países de ingresos bajos y medios, el cribado de la RD se basa principalmente en la evaluación del fondo de ojo realizada por oftalmólogos. Sin embargo, este enfoque resulta ineficiente frente a las limitaciones de recursos humanos y las barreras de acceso al sistema de salud previamente discutidas \parencites{piyasena2018}.

Según \cite{ortizbasso2023}, diferentes regiones del mundo han incorporado programas de telemedicina para la detección precoz de RD, logrando mejoras significativas en el acceso al sistema de salud, optimización de recursos y control de la enfermedad. Programas implementados en Inglaterra, Estados Unidos e India, han demostrado su efectividad para incrementar la cobertura y reducir complicaciones derivadas de la RD. En el contexto local, el estudio destaca que, tras la implementación de un programa de teleoftalmología en una zona rural de la provincia de La Pampa, la tasa anual de exámenes de fondo de ojos realizados en personas con diabetes aumentó significativamente, mostrando el potencial impacto positivo de estas iniciativas en áreas con recursos limitados.

Los avances en inteligencia artificial (IA) y, en particular, en aprendizaje profundo (DL), han permitido desarrollar algoritmos capaces de analizar imágenes de fondo de ojo con gran precisión. Estudios recientes han demostrado que la clasificación de fotografías de fondo de ojo con redes neuronales convolucionales (CNN), reconocidas por su capacidad para identificar patrones en imágenes, no solo logran detectar anomalías y patologías presentes en dichas imágenes, sino que también pueden establecer automáticamente el grado de RD en un paciente. Estos modelos, entrenados con diversos conjuntos de datos de imágenes oculares, han alcanzado métricas de rendimiento comparables a las de especialistas humanos, destacando su potencial en el diagnóstico automatizado de la RD \parencites{cuello2020}.

En 2018, la Administración de Drogas y Alimentos (FDA) de Estados Unidos, aprobó el primer dispositivo de detección autónoma con IA para diagnosticar RD, sin la necesidad de que un médico interprete la imagen o los resultados \parencites{abramoff2018}. Así mismo, en España, el Hospital Universitario Nuestra Señora de la Candelaria desarrolló su propia iniciativa para aplicar IA en el diagnóstico de la RD y, en 2020, firmó un acuerdo para obtener la certificación europea que permita su uso como dispositivo médico apto para la práctica clínica real \parencites{retina2020}.

En Argentina, proyectos como Retinar, desarrollados por instituciones nacionales como la UNICEN, el CONICET y el Hospital de Alta Complejidad El Cruce, han explorado soluciones innovadoras basadas en inteligencia artificial (IA). Entre los objetivos principales de este tipo de iniciativas se encuentra la implementación de plataformas tecnológicas que integren redes de captura de fotografías de fondo de ojo con oftalmólogos especialistas, para brindar diagnósticos e informes detallados allí donde no existe disponibilidad de profesionales. Estas soluciones utilizan dispositivos manejados por técnicos capacitados y asistencia por IA para garantizar la calidad de las capturas y realizar análisis preliminares de riesgo. Este enfoque busca optimizar el diagnóstico, reducir la carga de trabajo de los especialistas y mejorar la accesibilidad al tratamiento.
